% Generated draft: requires mathematical and source review.
\documentclass[11pt]{article}
\usepackage[margin=25mm]{geometry}
\usepackage{fontspec}
\setmainfont{Latin Modern Roman}
\usepackage{amsmath,amssymb,mathtools,unicode-math}
\setmathfont{Latin Modern Math}
\usepackage{enumitem,needspace,xurl,hyperref}
\setlength{\parindent}{0pt}
\setlength{\parskip}{5pt}
\setlength{\emergencystretch}{3em}
\newcommand{\N}{\mathbb{N}}
\newcommand{\Z}{\mathbb{Z}}
\newcommand{\Q}{\mathbb{Q}}
\newcommand{\R}{\mathbb{R}}
\newcommand{\C}{\mathbb{C}}
\newcommand{\F}{\mathbb{F}}
\newcommand{\A}{\mathbb{A}}
\newcommand{\PP}{\mathbb{P}}
\newcommand{\E}{\mathbb{E}}
\newcommand{\eps}{\varepsilon}
\newcommand{\dd}{\,\mathrm{d}}
\begin{document}
\section*{UnsolvedMath ID: 700319}
\textbf{Source problem code:} catalogue entry
\subsection{Definitions and mathematical statement}
A profinite group is a topological group isomorphic, as a topological group, to an inverse limit of finite groups endowed with the discrete topology. Say that a profinite group \(G\) satisfies the \emph{maximum condition for closed subgroups} if every ascending sequence of closed subgroups
\[
H_1\leq H_2\leq H_3\leq\cdots\leq G
\]
stabilizes: there exists \(n_0\in\N\), with \(n_0\geq 1\), such that \(H_n=H_{n_0}\) for every \(n\geq n_0\). Here closedness refers to the profinite topology, and the subgroups are not required to be normal. Equivalently, every nonempty family of closed subgroups of \(G\) has a member maximal under inclusion within that family.

\textbf{Problem (Yu. N. Mukhin, 1976; catalogue entry).} Characterize the profinite groups that satisfy the maximum condition for closed subgroups. No finiteness, countability, or metrizability hypothesis is imposed on \(G\).
\subsection{Short English statement}
Which profinite groups have the property that every increasing sequence of closed subgroups eventually becomes constant? The task is to characterize all such groups.
\subsection{Sources}
\begin{itemize}
\item Local UnsolvedMath registry: permanent ID \(700319\), exact source problem code \texttt{catalogue entry}, title \emph{Kourovka Notebook Problem 5.36}; supplied upstream revision \texttt{d371ee917ac9647fbb06cd6fc529deec6f3dda82}.
\item Supplied source attribution: \emph{The Kourovka Notebook (Unsolved Problems in Group Theory)}, problem 5.36, originally in the 5th issue (1976); proposed by Yu. N. Mukhin. The record identifies the 21st issue edition (2026), \texttt{arXiv:1401.0300v48}, as its source edition.
\item The supplied collection description names E. I. Khukhro and V. D. Mazurov as editors. The registry labels this problem \emph{open}; this is a supplied status label, not an independently verified mathematical or literature assessment.
\end{itemize}
\end{document}
